% Part I: Abstract, Introduction, Related Work

%% ═══════════════════════════════════════════════════════════════════════
%%  ABSTRACT
%% ═══════════════════════════════════════════════════════════════════════

\begin{abstract}
We introduce \emph{Judgment Geometry} (\JuGeo{}), a mathematical framework
for program verification in which programs are embedded as \emph{semantic sites}
(categories equipped with Grothendieck topologies), properties become
\emph{local sections} of presheaves over those sites, and verification reduces
to showing that \v{C}ech cohomology vanishes: $\Hcoh{1}(\Site,\mathcal{F})=0$.
This replaces the foundations of existing proof assistants---the Calculus of
Inductive Constructions (\LEAN{}, \Coq{}) and dependent refinement types
(\Fstar{})---with sheaf-theoretic descent, yielding a framework where local
specifications automatically glue into global guarantees by theorem rather than
by manual tactic construction.

The central algebraic object is the \emph{judgment 8-tuple}
$\J = \Jud{\coord}{\prop}{\carrier}{\evid}{\oblig}{\obstr}{\trust}{\prov}$,
which extends classical type-theoretic judgments with evidence bundles, a
seven-tier trust algebra $\Talg$, residual obligations, persistent
obstructions, and cryptographic provenance.  Over this algebra we build a
complete verification pipeline: specification satisfaction, equivalence
checking, bug detection via non-trivial $\Hcoh{1}$, obstruction-guided
program repair, monadic effect tracking, cover-guided program synthesis,
proof-carrying code with certificate chains, multi-agent treaty synthesis,
fragment-aware SMT dispatch, and a language of semantic moves that drives
proof search geometrically.

We present a working implementation (949 \Python{} files, 22 subsystems,
5\,653 classes) and evaluate on \expTotalPrograms{} programs across
\expDomainCount{} domains, achieving \expAccuracy{} verification accuracy
with \expPropsSum{} propositions verified in \expTimeMean{} mean time per
program.  \JuGeo{} generates proof-carrying \Python{} code and compiles
certificates to \LEAN{}.  We argue that Judgment Geometry opens a new field
at the intersection of algebraic geometry, formal methods, and AI-assisted
reasoning.
\end{abstract}

\section*{Keywords}
Sheaf theory, program verification, Grothendieck topology, \v{C}ech
cohomology, trust algebra, semantic site, proof-carrying code, Python

%% ═══════════════════════════════════════════════════════════════════════
\section{Introduction}\label{sec:intro}
%% ═══════════════════════════════════════════════════════════════════════

\subsection{The verification barrier}\label{sec:intro:barrier}

Formal verification promises mathematically certain software, yet the
promise remains unrealized for the vast majority of working programmers.
The obstacle is not theoretical---the metatheory of \LEAN{}, \Coq{}, and
\Fstar{} is impeccable---but practical.  To verify a 200-line \Python{}
sorting algorithm in \LEAN{}, one must rewrite it in Lean's functional
language, formulate inductive invariants, and construct tactic proofs
whose combined length routinely exceeds the program itself.  \Fstar{}
reduces the tactic burden by encoding specifications as dependent
refinement types, but demands that the programmer annotate every function
with computation types $\mathsf{M}\;a\;\mathit{wp}$ and reason about
weakest-precondition transformers.  \Dafny{} lowers the entry bar with
Hoare-style annotations, yet still requires the target program to be
rewritten in Dafny's own language and the user to supply loop invariants
by hand.

In every case, verification imposes a \emph{language barrier} (rewrite in
the host language), a \emph{proof barrier} (construct explicit evidence),
and a \emph{composition barrier} (manually compose local results into global
ones).  These three barriers conspire to keep formal verification confined to
safety-critical niches: avionics, cryptographic libraries, compiler backends.
The rest of the software world---web services, data pipelines, machine-learning
infrastructure---remains unverified.

\subsection{The geometric insight}\label{sec:intro:insight}

The central observation of this paper is that programs \emph{already have
geometric structure}, and that verification is \emph{already a descent
problem} in the sense of Grothendieck's algebraic geometry.

Consider a Python module with functions $f_1, \ldots, f_k$.  Each function
can be viewed as a \emph{chart} on a manifold: it provides a local coordinate
system in which a portion of the program's behavior is described.  The
interfaces between functions---shared state, call-return contracts,
exception propagation---are the \emph{overlaps} between charts.  A
specification (``the output list is sorted,'' ``the balance is
non-negative'') is a \emph{local section}: it holds on each chart separately.
The question ``does the specification hold globally?'' is exactly the
question of whether local sections \emph{glue}---that is, whether they
satisfy the descent condition of a sheaf.

This is not an analogy.  It is a precise mathematical correspondence
that we formalize in full.
A program~$P$ determines a small category $\Site_P$ whose objects are
\emph{coordinates}---modules, classes, functions, branches, loop bodies,
expressions---and whose morphisms are typed arrows: restrictions (function
$\to$ body), inclusions (branch $\to$ parent conditional), transports
(call site $\to$ callee), and refinements (coordinate $\to$ finer
decomposition).  We equip this category with a Grothendieck topology
$\Cov$: the covering families are the collections of coordinates that
``jointly observe everything'' about a region.  For instance, the two
branches of an \texttt{if/else} form a covering family of the
conditional, and the functions of a module form a covering family of the
module coordinate.  The triple of axioms---identity, stability, and
transitivity---is verified constructively (\S\ref{sec:sites}).

Properties of the program are modeled as presheaves
$\mathcal{F} : \Site_P^{\mathrm{op}} \to \mathbf{Set}$: for each
coordinate~$U$, the set $\mathcal{F}(U)$ contains the verification
conditions local to~$U$, and for each morphism $f : V \to U$ the
restriction $\res_f : \mathcal{F}(U) \to \mathcal{F}(V)$ specializes a
global condition to a sub-region.  The sheaf condition
\[
  \mathcal{F}(U) \;\xrightarrow{\;\sim\;}\;
  \varprojlim\!\Bigl(
    \prod_i \mathcal{F}(U_i) \;\rightrightarrows\;
    \prod_{i,j} \mathcal{F}(\overlap{U_i}{U_j})
  \Bigr)
\]
is exactly the statement that locally verified properties lift to a global
guarantee.  Verification therefore reduces to a cohomological computation:

\begin{theorem}[Descent criterion, informal]\label{thm:descent-informal}
Let $\Site_P$ be the semantic site of a program~$P$, and let $\mathcal{F}$
be the presheaf of verification conditions.  The specification is globally
satisfied if and only if the first \v{C}ech cohomology vanishes:
\[
  \Hcoh{1}(\Site_P,\, \mathcal{F}) \;=\; 0.
\]
When $\Hcoh{1} \neq 0$, its non-trivial classes are \emph{obstructions}---they
identify exactly which overlaps between program regions fail to reconcile, and
they guide automated repair.
\end{theorem}

A formal statement and proof appear in \S\ref{sec:descent}; the informal
version suffices to convey the architecture.

\begin{example}[Sorting verification]\label{ex:sort-intro}
Consider a \Python{} function \texttt{merge\_sort(xs)} that splits
\texttt{xs} into halves, recursively sorts each half, and merges the
results.  The semantic site has coordinates for the split, the two
recursive calls, and the merge.  The specification ``output is sorted
and a permutation of input'' is a presheaf section.  On each coordinate
the section is verified locally (the split preserves elements; each
recursive call returns a sorted list; the merge of two sorted lists is
sorted).  The overlaps---shared variables between split and recursion,
between recursion and merge---form the \v{C}ech $1$-cochains.  The
coboundary computes restriction differences: if the recursive call's
postcondition matches the merge's precondition, the $1$-cocycle is a
coboundary and $\Hcoh{1} = 0$.  The specification holds globally.
\end{example}

\subsection{The JuGeo framework}\label{sec:intro:framework}

We build on this geometric insight to construct a complete verification
framework.  The architecture has three layers.

\paragraph{Layer 1: Semantic sites.}
Given a \Python{} program~$P$, the \JuGeo{} pipeline parses the AST and
constructs a site $\Site_P = (\mathcal{C}_P, \Cov_P)$ whose coordinates
are the syntactic and semantic regions of~$P$ and whose covering families
encode how control flow, scoping, and module structure decompose~$P$ into
overlapping regions.  Coordinates are typed by kind---\texttt{MODULE},
\texttt{FUNCTION}, \texttt{INTERFACE}, \texttt{REGION}, \texttt{THEOREM},
\texttt{TEST}---and morphisms carry semantic labels (restriction, inclusion,
transport, refinement).  The topology satisfies the three Grothendieck axioms:
identity covers, stability under pullback, and transitivity.  Site complexity
scales linearly in AST size (\S\ref{sec:eval}).

\paragraph{Layer 2: Judgment algebra.}
Over the site we define judgments.  Unlike classical type-theoretic
judgments $\Gamma \vdash a : A$ (three components), a \JuGeo{} judgment
is an 8-tuple:
\begin{equation}\label{eq:judgment}
  \J \;=\; \Jud{\coord}{\prop}{\carrier}{\evid}{\oblig}{\obstr}{\trust}{\prov}
\end{equation}

\begin{definition}[Judgment 8-tuple, informal]\label{def:judgment-intro}
The components of~\eqref{eq:judgment} are:
\begin{enumerate}[label=(\roman*),itemsep=1pt]
  \item $\coord \in \Ob(\Site_P)$: the \emph{coordinate}---the program
        region to which the judgment applies.
  \item $\prop$: the \emph{proposition}---a claim about behavior, typed as
        structural, behavioral, relational, resource, or semantic.
  \item $\carrier$: the \emph{carrier type}---what the claim is about,
        possibly dependent on $\coord$.
  \item $\evid$: the \emph{evidence bundle}---a multiset of evidence items,
        each tagged with a channel (solver, runtime, oracle, human) and a
        trust level.
  \item $\oblig$: \emph{residual obligations}---verification conditions not
        yet discharged.
  \item $\obstr$: \emph{obstructions}---persistent blocking records;
        first-class cohomology classes, never silently erased.
  \item $\trust \in \Talg$: the \emph{trust annotation}---an element of
        the trust ordered algebra (not a scalar but an algebraic object
        with six operations).
  \item $\prov$: the \emph{provenance}---a record of the derivation
        history (solver, runtime, oracle, human, or composed).
\end{enumerate}
\end{definition}

The trust algebra $\Talg = (\Eadm, \tleq, \tjoin, \tatten, \tpromote,
\tdemote)$ has seven tiers:
\[
  \Tcontra \;\tleq\; \Tunver \;\tleq\; \Tcopilot \;\tleq\;
  \Toracle \;\tleq\; \Truntime \;\tleq\; \Tsolver \;\tleq\; \Tproof
\]
with three invariants: \emph{no silent promotion} (every upward move
requires explicit justification in an append-only audit log),
\emph{conservative join} ($\trust_1 \tjoin \trust_2 = \min(\trust_1,
\trust_2)$), and \emph{challenge conservativity} (a challenge can only
lower trust).  We write $\trust_1 \tleq \trust_2$ for the partial order
and note that this is a bounded distributive lattice with $\Tcontra$ as
bottom and $\Tproof$ as top.

Five algebraic operations---restriction $\res_{U \to V}$, transport
$\tau_{U \to V}$, composition $\J_1 \circ \J_2$, comparison, and
join---satisfy functoriality and coherence laws, and the classical
structural rules (weakening, contraction, exchange, cut) are admissible
(\S\ref{sec:algebra}).

\paragraph{Layer 3: Cohomological descent.}
We form the \v{C}ech complex
\[
  \Cech^0(\mathfrak{U}, \mathcal{F})
  \;\xrightarrow{\;\delta^0\;}\;
  \Cech^1(\mathfrak{U}, \mathcal{F})
  \;\xrightarrow{\;\delta^1\;}\;
  \Cech^2(\mathfrak{U}, \mathcal{F})
  \;\to\; \cdots
\]
where $\mathfrak{U} = \{U_i\}$ is a covering family in $\Site_P$ and
$\mathcal{F}$ is the presheaf of judgments.  The $0$-cochains are local
judgments on each~$U_i$; the $1$-cochains record mismatches on overlaps
$\overlap{U_i}{U_j}$; the coboundary~$\delta^0$ computes restriction
differences.  The cohomology class $\Hcoh{1}(\Site_P, \mathcal{F})$
classifies obstructions to gluing: it vanishes if and only if local
judgments are globally consistent.  When it does not vanish, the
non-trivial cocycles localize the failure to specific overlaps and
guide the repair engine (\S\ref{sec:repair}).

\subsection{Capabilities}\label{sec:intro:capabilities}

This geometric architecture unifies ten verification tasks within a single
framework:

\begin{enumerate}[label=(\arabic*),leftmargin=2em,itemsep=2pt]
  \item \textbf{Specification satisfaction}: verify that a program meets its
        specification by computing $\Hcoh{1}(\Site_P, \mathcal{F}_{\mathrm{spec}}) = 0$.
  \item \textbf{Equivalence checking}: prove two programs $P, Q$ extensionally
        equivalent by constructing an isomorphism of presheaves
        $\mathcal{F}_P \cong \mathcal{F}_Q$ over a common refinement site.
  \item \textbf{Bug detection}: identify bugs as non-trivial classes in
        $\Hcoh{1}$; each obstruction localizes the bug to a specific overlap.
  \item \textbf{Program repair}: use obstructions as repair directives---each
        cocycle suggests a minimal code patch that would kill the cohomology class.
  \item \textbf{Effect tracking}: model monadic effects (\texttt{IO},
        \texttt{State}, exceptions, async) as presheaves; the descent condition
        ensures effect consistency across module boundaries.
  \item \textbf{Program synthesis}: generate programs by constructing covers
        that admit global sections---the cover topology guides synthesis toward
        well-structured code.
  \item \textbf{Proof-carrying code}: attach certificate chains to programs;
        each certificate is a verified judgment with trust annotation and
        provenance, checkable independently of the prover.
  \item \textbf{Treaty synthesis}: reconcile multi-agent verification
        (solver + oracle + runtime) via treaty objects that negotiate
        consistent trust-annotated judgments across agents.
  \item \textbf{SMT dispatch}: route verification conditions to appropriate
        SMT fragments (\QFLIA{}, \QFLRA{}, \QFBV{}, \QFUF{}) via
        fragment-aware theory classification.
  \item \textbf{Semantic moves}: drive proof search via a vocabulary of
        geometric operations---\VERIFY{}, \CONSTRUCT{}, \NORMALIZE{},
        \GLUE{}, \RESTRICT{}, \TRANSPORT{}, \REPAIR{}, \PROMOTE{},
        \CHALLENGE{}, \ESCALATE{}---each with well-defined pre- and
        postconditions on the judgment lattice.
\end{enumerate}

\subsection{Experimental summary}\label{sec:intro:experiments}

We evaluate \JuGeo{} on \expTotalPrograms{} \Python{} programs spanning
\expDomainCount{} domains (sorting, data structures, mathematical
algorithms, string processing, state machines, web handlers).  The
framework verifies all \expVerified{} programs with \expAccuracy{}
accuracy, discharging \expPropsSum{} propositions with zero obstructions
($\Hcoh{1} = 0$ on all \expHOneZero{} programs).  Mean verification
time is \expTimeMean{} per program (range: \expTimeMin{}--\expTimeMax{}).
Per-domain statistics confirm uniformity: coordinates range from
\expCoordsMin{} to \expCoordsMax{} (mean \expCoordsMean{}), propositions
from \expPropsMin{} to \expPropsMax{} (mean \expPropsMean{}).

The implementation comprises 949 \Python{} files organized into 22
subsystems with 5\,653 classes.  It generates proof-carrying \Python{}
code with embedded certificates and compiles trust-annotated judgments to
\LEAN{} for independent verification.

\subsection{Contributions}\label{sec:intro:contributions}

\begin{enumerate}
  \item A mathematical framework---\emph{Judgment Geometry}---that replaces
        CIC and dependent types with sheaf-theoretic descent over semantic
        sites as the foundation for program verification (\S\ref{sec:sites}--\S\ref{sec:descent}).
  \item The judgment 8-tuple $\J$ with trust algebra $\Talg$, extending
        classical proof-theoretic judgments with evidence provenance, trust
        gradation, residual obligations, and persistent obstructions
        (\S\ref{sec:algebra}).
  \item A cohomological diagnostic: $\Hcoh{1} = 0$ certifies correctness;
        $\Hcoh{1} \neq 0$ classifies and localizes failures
        (\S\ref{sec:descent}).
  \item A complete implementation and evaluation demonstrating
        \expAccuracy{} accuracy on \expTotalPrograms{} programs across
        \expDomainCount{} domains (\S\ref{sec:eval}).
  \item Evidence that Judgment Geometry opens a new field at the
        intersection of algebraic geometry, formal methods, and AI-assisted
        reasoning (\S\ref{sec:future}).
\end{enumerate}

\subsection{Paper organization}\label{sec:intro:organization}

The remainder of this paper is organized as follows.
\S\ref{sec:related} surveys related work.
\S\ref{sec:sites} develops semantic sites and their Grothendieck topologies.
\S\ref{sec:algebra} defines the judgment 8-tuple and its algebraic
operations.
\S\ref{sec:descent} presents the \v{C}ech cohomology framework and the
descent criterion.
\S\ref{sec:trust} develops the trust algebra.
\S\ref{sec:pipeline} describes the SMT dispatch, semantic moves, effect
tracking, treaty synthesis, and proof-carrying code subsystems.
\S\ref{sec:repair} presents obstruction-guided program repair.
\S\ref{sec:eval} reports experimental results.
\S\ref{sec:future} discusses future directions and open problems.

%% ═══════════════════════════════════════════════════════════════════════
\section{Related Work}\label{sec:related}
%% ═══════════════════════════════════════════════════════════════════════

We survey the landscape of program verification and position \JuGeo{}
relative to each major paradigm.  A summary comparison appears in
Table~\ref{tab:comparison}.

\subsection{Type-theoretic proof assistants}\label{sec:related:type-theory}

\paragraph{Lean 4.}
\LEAN{}~\cite{lean4} is founded on the Calculus of Inductive Constructions
(CIC), a dependent type theory with inductive families, universe
polymorphism, and a small trusted kernel.  Verification proceeds by
constructing proof terms (via tactic scripts) that the kernel type-checks.
\LEAN{}'s strengths are well understood: a rich mathematical library
(\texttt{Mathlib}), strong metaprogramming via monadic \texttt{do}-notation,
and a growing community.  Its limitations for software verification are
equally clear.  The target program must be rewritten in Lean's functional
language; there is no mechanism for verifying unmodified Python, Java, or
C code.  Proofs are manual---the user must supply induction hypotheses,
case splits, and rewriting chains.  Perhaps most critically, \LEAN{} makes
no distinction between a proof found by a Fields Medalist in six months and
a proof found by an LLM in six seconds: both receive the same kernel
\texttt{Qed}.  There is no trust gradation, no provenance tracking, and no
notion of partial verification.

\JuGeo{} differs in foundation (sheaf-theoretic descent vs.\ CIC),
target language (native \Python{} vs.\ Lean), proof style (automatic
cohomological descent vs.\ manual tactics), and trust model (seven-tier
algebra with audit trail vs.\ binary accept/reject).  Where \LEAN{}
verifies by type-checking a proof term, \JuGeo{} verifies by computing
$\Hcoh{1} = 0$.  Where \LEAN{} reports ``type error'' on failure, \JuGeo{}
returns a structured obstruction localizing the failure to specific
overlaps between program regions.

\paragraph{Coq.}
\Coq{}~\cite{coq} shares \LEAN{}'s CIC foundation but uses the Gallina
specification language with the Ltac tactic engine.  Code verified in
\Coq{} must be written in Gallina or extracted to OCaml via the extraction
mechanism---a one-directional pipeline that cannot verify existing OCaml
programs.  \Coq{}'s proof terms are certifiable (the \texttt{Print
Assumptions} command lists axioms used), but there is no mechanism for
tracking which solver or oracle contributed to a proof, no trust algebra,
and no first-class obstruction theory.

\JuGeo{} replaces the extraction pipeline with native \Python{} analysis:
programs are verified in place, and certificates are embedded directly in
the source.  The judgment 8-tuple subsumes Gallina's three-component
judgment $\Gamma \vdash t : T$ as a degenerate case
(see \S\ref{sec:algebra}).

\paragraph{F*.}
\Fstar{}~\cite{fstar} extends dependent types with an effect system,
encoding specifications as weakest-precondition (WP) transformers in
computation types $\mathsf{M}\;a\;\mathit{wp}$.  Verification conditions
are discharged automatically via Z3, yielding a higher degree of automation
than \LEAN{} or \Coq{}.  However, \Fstar{} requires the target program to
be written in \Fstar{}'s own syntax with explicit effect annotations, and
its WP calculus---while powerful---does not model the \emph{spatial
structure} of the program under verification.  Two functions verified
independently in \Fstar{} are composed via monadic bind; there is no
topological account of \emph{which overlaps} were checked or \emph{why}
composition is sound.

\JuGeo{} replaces WP transformers with presheaf sections and monadic
composition with sheaf descent.  The trust algebra replaces \Fstar{}'s
binary verified/unverified distinction with a seven-tier lattice that
distinguishes solver proofs from runtime witnesses from oracle proposals.
Properties are inferred from the program's AST rather than annotated by
the programmer.

\subsection{Annotation-based verifiers}\label{sec:related:annotation}

\paragraph{Dafny.}
\Dafny{}~\cite{dafny} compiles method-level Hoare triples into the
Boogie intermediate verification language, which discharges conditions
via Z3.  \Dafny{}'s strength is accessibility: pre/postconditions and loop
invariants are written in a familiar imperative syntax.  Its weakness is
diagnostic poverty---when verification fails, the user receives a
counterexample from Z3 but no structured explanation of \emph{why} the
proof attempt failed or \emph{where} in the program's modular structure
the inconsistency lives.

\JuGeo{} provides richer geometric diagnostics via its obstruction theory.
A \Dafny{} counterexample is a zero-dimensional object (a single failing
execution trace); a \JuGeo{} obstruction is a one-dimensional cohomology
class that identifies the \emph{overlap} between two program regions where
specifications disagree.  Formally, let $U_i, U_j$ be two coordinates
(e.g., a function and its caller) with local judgments $\J_i, \J_j$.  A
\Dafny{} counterexample is a point $x \in \overlap{U_i}{U_j}$ witnessing
$\J_i|_{U_i \cap U_j} \neq \J_j|_{U_i \cap U_j}$; a \JuGeo{}
obstruction is the \emph{class} $[\alpha] \in \Hcoh{1}(\Site_P,
\mathcal{F})$ represented by the cocycle $\alpha_{ij} = \res_i(\J_i) -
\res_j(\J_j)$.  This higher-dimensional structure enables
obstruction-guided repair (\S\ref{sec:repair}), which has no analogue in
\Dafny{}.

\paragraph{Liquid Haskell.}
Liquid Haskell~\cite{liquidhaskell} adds refinement types to Haskell,
enabling specification of properties like sortedness or positivity as
type-level predicates discharged by an SMT solver.  It operates on Haskell
source and inherits the functional paradigm's strengths (purity,
algebraic data types, pattern matching) and weaknesses (no imperative
code, no mutation, no native \Python{} support).

\JuGeo{} targets imperative \Python{} with mutation, side effects, and
dynamic typing---a fundamentally different verification challenge.  Where
Liquid Haskell encodes specifications as refinement types, \JuGeo{} encodes
them as presheaf sections.  The geometric approach naturally handles
imperative features (mutable state as restriction morphisms, exception
flow as transport morphisms) that have no direct type-theoretic counterpart.

\subsection{Model checking and symbolic execution}\label{sec:related:model-checking}

\paragraph{CBMC and KLEE.}
Bounded model checkers such as CBMC~\cite{cbmc} and symbolic execution
engines such as KLEE~\cite{klee} explore program state spaces up to a
fixed depth~$k$, encoding reachability as SAT/SMT queries.  They excel at
finding concrete bugs in C code but provide no \emph{unbounded}
guarantees: a program verified to depth~$k$ may fail at depth~$k{+}1$.
Their output is a single counterexample trace---again a zero-dimensional
object---with no cohomological structure.

\JuGeo{}'s cohomological descent provides unbounded verification: the
vanishing $\Hcoh{1}(\Site_P, \mathcal{F}) = 0$ certifies the
specification for \emph{all} inputs, not merely those reachable within
a bound.  Moreover, every judgment carries a trust annotation---a
bounded model-checking result would enter the \JuGeo{} pipeline at
the $\Truntime$ tier, strictly below $\Tsolver$ and $\Tproof$,
with the depth bound recorded in provenance.  This enables \JuGeo{} to
\emph{integrate} BMC results into a broader verification without
conflating bounded evidence with unbounded proof.

\subsection{Sheaf theory in computer science}\label{sec:related:sheaf-cs}

Abramsky and Brandenburger~\cite{abramsky2011} introduced
sheaf-theoretic models of contextuality in quantum mechanics, showing
that non-locality corresponds to the failure of a presheaf to be a sheaf.
This work demonstrated that sheaf cohomology has computational content
beyond pure mathematics.  Goguen~\cite{goguen1992} used sheaves to model
information systems, and more recently, Robinson~\cite{robinson2014}
applied sheaves to sensor fusion and data integration.
Curry~\cite{curry2014} developed sheaves on cell complexes for topological
signal processing.  In formal methods, Pavlovic and
colleagues~\cite{pavlovic2012} explored categorical models of security
protocols with sheaf-like gluing conditions.

\JuGeo{} is, to our knowledge, the first system to build a \emph{complete
end-to-end verification pipeline} from sheaf theory---from AST parsing to
certificate generation to \LEAN{} compilation.  Prior work used
sheaves as \emph{models} of existing phenomena (contextuality, information
flow); \JuGeo{} uses sheaves as the \emph{computational mechanism} for
verification itself.  The key innovations---the judgment 8-tuple, the trust
algebra, the semantic site construction for arbitrary programs, the
obstruction-guided repair loop, and the treaty synthesis protocol---have
no precedent in the sheaf-theory literature.

The closest technical antecedent is Abramsky's observation that the failure
of a presheaf to satisfy the sheaf condition has computational meaning.
\JuGeo{} operationalizes this observation: the failure of the presheaf of
verification conditions to be a sheaf is precisely the presence of bugs,
and the cohomology class of that failure is precisely the bug's
localization.

\subsection{AI-assisted proof}\label{sec:related:ai}

The use of large language models (LLMs) for proof generation has
accelerated rapidly.  AlphaProof~\cite{alphaproof2024} combines a
language model with AlphaZero-style reinforcement learning to solve
International Mathematical Olympiad problems in \LEAN{}.
LeanDojo~\cite{leandojo2023} provides a benchmark and training pipeline
for LLM-based tactic prediction.  Draft-Sketch-Prove~\cite{dsp2022}
uses informal proofs as sketches for formal proof generation.

These systems use AI as the \emph{primary proof mechanism}: the LLM
generates tactic steps or proof terms, and the kernel checks them.  The
risk is well known---a model that generates correct-looking proofs 95\%
of the time provides no formal guarantee on the remaining 5\%.

\JuGeo{} takes a fundamentally different approach.  AI (specifically,
GitHub Copilot and GPT-4 class models) serves as an \emph{oracle tier}
within the trust algebra.  An oracle-generated claim enters the system at
trust level $\Tcopilot$ or $\Toracle$---strictly below $\Tsolver$ and
$\Tproof$---and cannot be silently promoted.  The oracle's contribution
is recorded in the provenance field~$\prov$, subject to challenge via the
$\CHALLENGE{}$ semantic move, and attenuated by the $\tatten$ operator
if contradictory evidence is found.  AI assists verification but does not
replace it: every oracle claim must either be independently confirmed
(promoting it to a higher tier with explicit justification logged in the
audit trail) or remain at its ceiling, clearly marked in the certificate.

This design reflects a deeper philosophical commitment: in \JuGeo{}, the
geometry of the proof---the site, the covering, the descent condition---is
the source of truth, not the oracle.  An LLM can \emph{suggest} that a
function satisfies its postcondition, but only the vanishing of
$\Hcoh{1}$ \emph{proves} it.

\subsection{Comparison summary}\label{sec:related:summary}

Table~\ref{tab:comparison} summarizes the comparison across eight
dimensions.  \JuGeo{} is unique in its sheaf-theoretic foundation,
automatic proof style, seven-tier trust tracking, cohomological bug
detection, obstruction-guided repair, and multi-agent treaty synthesis.

The comparison reveals a structural gap in the verification landscape.
Existing tools occupy two poles: \emph{deep but narrow} (type-theoretic
proof assistants that require language rewriting and manual proofs) and
\emph{broad but shallow} (linters and model checkers that scale easily
but provide limited guarantees).  \JuGeo{} occupies a third position:
deep guarantees via cohomological descent, applied broadly to native
\Python{} code with automatic proof search.  The trust algebra provides
the mechanism for bridging the poles---a \LEAN{} kernel proof and a
runtime test can coexist in the same judgment, each at its appropriate
trust tier, composed via the conservative join $\tjoin$.

\begin{table}[t]
\centering
\caption{Comparison of \JuGeo{} with existing verification frameworks.}
\label{tab:comparison}
\small
\begin{tabularx}{\textwidth}{l *{5}{>{\centering\arraybackslash}X}}
\toprule
\textbf{Feature} & \textbf{\LEAN{}} & \textbf{\Fstar{}} & \textbf{\Coq{}}
  & \textbf{\Dafny{}} & \textbf{\JuGeo{}} \\
\midrule
Foundation
  & CIC & Dep.\ Types & CIC & Hoare Logic & Sheaf Descent \\[3pt]
Target Language
  & Lean & \Fstar{} & Gallina & Dafny & \Python{} \\[3pt]
Proof Style
  & Tactics & Refinement & Tactics & Annotations & Automatic \\[3pt]
Trust Tracking
  & No & No & No & No & Yes (7 tiers) \\[3pt]
Cohomological
  & No & No & No & No & Yes ($\Hcoh{1}$) \\[3pt]
Bug Detection
  & No & No & No & Counterex. & Obstruction \\[3pt]
Program Repair
  & No & No & No & No & Obstr.-guided \\[3pt]
Multi-agent
  & No & No & No & No & Treaty synth. \\
\bottomrule
\end{tabularx}
\end{table}

\begin{remark}[Scope of comparison]\label{rem:scope}
The ``No'' entries in Table~\ref{tab:comparison} indicate that the
feature is not a built-in, first-class capability of the system.
Third-party tools or plugins may provide partial support---for instance,
\LEAN{}'s \texttt{Mathlib} tactics automate some proofs, and \Dafny{}'s
counterexamples provide limited diagnostic information.  The comparison
reflects the core framework, not the ecosystem.
\end{remark}

\begin{remark}[Complementarity]\label{rem:complementarity}
\JuGeo{} does not \emph{replace} \LEAN{}, \Fstar{}, or \Coq{}; it
\emph{complements} them.  \JuGeo{}'s proof-carrying code subsystem
compiles trust-annotated judgments to \LEAN{} for independent
verification, and the trust algebra can ingest \LEAN{} kernel
\texttt{Qed} as evidence at the $\Tproof$ tier.  The relationship
is symbiotic: \JuGeo{} provides the geometric decomposition and
automation layer; \LEAN{} provides the trusted kernel for the
highest-assurance tier.  We envision a workflow in which \JuGeo{}
handles the 90\% of verification conditions that can be discharged
automatically, and the remaining 10\%---the deep mathematical
lemmas---are escalated to \LEAN{} via the $\ESCALATE{}$ semantic move,
with the resulting proofs ingested back at the $\Tproof$ tier.
\end{remark}
